% Original applied explanation; see Chapter04/SOURCES.md.
\subsubsection{ছবির প্রদর্শনের মাপ ও ফাইলের মাপ আলাদা}
ক্যাম্পাসের বড় ছবির \texttt{width} কমালে পর্দায় ছোট দেখা যায়; তাতে মূল ফাইলের byte-সংখ্যা নিজে থেকে কমে না। ধীর সংযোগে দ্রুত দেখাতে উপযুক্ত মাপের ছবির ফাইল তৈরি করতে হয়। ছবির প্রস্থ ও উচ্চতার অনুপাত না মিলিয়ে দুটোই বদলালে ছবি বিকৃত হতে পারে।

ছবি দেখা না গেলে \texttt{src}-এর ফোল্ডার, নাম, বড়/ছোট হাতের অক্ষর এবং ফাইলের ধরন মিলিয়ে দেখবে। তথ্যবহুল ছবির \texttt{alt} তার প্রয়োজনীয় অর্থ দেয়; কেবল সাজসজ্জার ছবি হলে খালি \texttt{alt=""} ব্যবহার করা যায়।
